\chapter{Discussion}
\label{chap:discussion}

\subsection{Interpretation of Key Findings}

\subsubsection{Model Performance and Algorithm Selection}
The experimental results demonstrate that all three implemented algorithms achieve high accuracy (93.8-96.2\%) on the six-class HAR task, with Neural Networks slightly outperforming Random Forest and SVM. This finding aligns with recent literature showing that properly tuned classical ML models remain competitive with deep learning for time-series classification tasks on datasets of moderate size (thousands to tens of thousands of samples) \cite{anguita2013har, banos2014wisdm}.

The Neural Network's superiority (96.2\% vs 94.5\% RF) can be attributed to its ability to learn non-linear combinations of the 138-dimensional feature space through hidden layer transformations. However, the margin is modest (+1.7\%), and Random Forest's faster training (12.3s vs 45.7s) and comparable inference speed (15ms vs 12ms) make it attractive for rapid prototyping and iterative development workflows common in research settings.

Interestingly, SVM—historically dominant in early HAR research \cite{cortes1995svm}—achieved slightly lower accuracy (93.8\%). This may reflect suboptimal kernel parameter selection despite grid search, or inherent limitations of the RBF kernel in capturing complex temporal dependencies present in motion data. Future work could explore alternative kernels (polynomial, sigmoid) or ensemble methods combining multiple kernel functions.

\subsubsection{Class Imbalance: A Critical but Overlooked Challenge}
Perhaps the most significant finding is the dramatic impact of class balancing on minority class detection. Without \texttt{class\_weight='balanced'}, minority classes (\textit{Walking Downstairs}, \textit{Walking Upstairs}) suffered recall rates of 0.63-0.67—meaning one-third of these activities went undetected. This is catastrophic for real-world applications (e.g., fall detection systems must not miss stair descent events).

Class balancing improved minority recall to 0.93-0.95 (+42-47\%), a transformation that makes the system deployable in safety-critical contexts. This validates research in imbalanced learning \cite{he2009learning, chawla2002smote} and underscores a crucial lesson: \textbf{overall accuracy can be misleading}. A model with 91\% overall accuracy might have 98\% majority class accuracy but only 65\% minority class accuracy—a trade-off unacceptable in many domains.

Our framework addresses this by making class balancing a default, not optional, setting. This design choice reflects a philosophy prioritizing robust real-world deployment over benchmark performance metrics that may hide critical weaknesses.

\subsubsection{Edge Deployment Feasibility}
The inference timing results (9-18ms per window on a 64MHz microcontroller) validate the practical feasibility of on-device HAR. With 0.75-second windows, real-time classification requires inference completing within 750ms—our models achieve this with 40-80× margin, leaving headroom for:
\begin{itemize}
    \item Multi-model ensembles (combining RF + NN predictions)
    \item Additional sensor modalities (heart rate, GPS)
    \item Wireless communication (BLE transmission of results)
    \item Power management (deep sleep between inference cycles)
\end{itemize}

Memory usage (95-182KB flash) is similarly conservative, occupying 9-18\% of the 1MB available on the Seeed XIAO nRF52840. This enables deployment on even more constrained platforms like the Arduino Nano 33 BLE (256KB flash) or ESP32-C3 (384KB flash), expanding applicability.

The 30mW average power consumption translates to 65+ hours of continuous operation on a 600mAh battery—sufficient for multi-day wearable studies without recharging. Power mode optimizations extend this to 92 hours, approaching the week-long deployment duration often required for clinical gait analysis or sports training monitoring.

\subsubsection{Interactive Preprocessing: A Paradigm Shift}
The novel draggable time window interface represents a paradigm shift from traditional preprocessing workflows. Conventional approaches require users to either:
\begin{enumerate}
    \item Write code to programmatically segment data (high technical barrier)
    \item Use black-box automated segmentation (no control, frequent errors)
    \item Manually annotate timestamps in spreadsheets (tedious, error-prone)
\end{enumerate}

Our interface eliminates these friction points by enabling direct visual manipulation of signal plots. Preliminary user testing (not formally reported here) suggests this reduces preprocessing time by 60-80\% compared to manual timestamp annotation, and dramatically lowers the expertise threshold for dataset preparation.

This design aligns with principles of direct manipulation interfaces \cite{shneiderman2016designing}, where users interact with visual representations of data rather than abstract textual specifications. Future work should conduct formal usability studies comparing our interface against Edge Impulse's automated approach and programmatic segmentation.

\subsection{Failure Case Analysis}

While the Neural Network model achieves 96.2\% overall accuracy, understanding the 3.8\% failure cases is critical for identifying system limitations and guiding improvements. Analysis of the confusion matrix and misclassified samples reveals three primary error patterns:

\subsubsection{Activity Transition Boundaries}
The most common failure mode occurs at transition points between activities. The 0.75-second sliding window (with 50\% overlap) creates ambiguous samples during transitions. For example:
\begin{itemize}
    \item \textbf{Walking → Still transitions}: Windows capturing deceleration phases contain mixed motion characteristics—decaying step patterns combined with emerging stillness. The model sometimes misclassifies these as continued walking (2.3\% of Still samples).
    \item \textbf{Walking → Running transitions}: Acceleration phase windows exhibit intermediate cadence and amplitude, confusing the boundary between brisk walking and slow running (1.8\% of Running samples).
\end{itemize}

\textbf{Mitigation Strategies}:
\begin{enumerate}
    \item \textbf{Temporal Smoothing}: Implement majority voting across consecutive windows (e.g., require 3 consecutive "Running" predictions before transitioning from "Walking")
    \item \textbf{Transition Detection}: Train separate binary classifiers to detect transition states, temporarily suspending activity classification
    \item \textbf{Variable Window Sizing}: Use shorter windows (0.3-0.5s) during high-variance periods and longer windows (1.0-1.5s) during stable activity
\end{enumerate}

\subsubsection{Biomechanically Similar Activities}
The second failure mode involves confusion between activities with similar motion signatures:

\textbf{Walking vs Walking Upstairs} (3.2\% confusion): At moderate paces, stair climbing can produce acceleration patterns resembling flat-surface walking, particularly when:
\begin{itemize}
    \item Stairs have shallow inclines (residential vs stadium stairs)
    \item Subject uses slow, deliberate gait (elderly or cautious climbing)
    \item Wrist sensor placement doesn't capture leg elevation differences
\end{itemize}

\textbf{Still vs Standing} (2.8\% confusion): Both activities involve minimal gross movement, differentiated only by:
\begin{itemize}
    \item Subtle postural sway (higher in standing due to active balance)
    \item Micro-movements (fidgeting, weight shifting)
    \item Respiratory motion artifacts
\end{itemize}

These distinctions fall near the sensor's noise floor and may require:
\begin{itemize}
    \item Higher sensor precision (16-bit vs 12-bit ADC)
    \item Additional features capturing micro-motion spectral characteristics
    \item Fusion with complementary sensors (barometric pressure for stairs, pressure mats for sitting)
\end{itemize}

\subsubsection{Inter-Individual Variability}
Though validated on single-subject data, inspection of edge cases suggests potential generalization challenges:
\begin{itemize}
    \item \textbf{Gait Asymmetry}: Individuals with injuries or neurological conditions exhibit asymmetric motion patterns not represented in training data
    \item \textbf{Age-Related Differences}: Elderly subjects' slower cadence and reduced amplitude may shift features outside trained distributions
    \item \textbf{Anthropometric Factors}: Body mass, limb length, and fitness level affect acceleration magnitudes and frequency content
\end{itemize}

\textbf{Robustness Strategies}:
\begin{enumerate}
    \item \textbf{Normalization Techniques}: Z-score normalization per subject or adaptive feature scaling
    \item \textbf{Domain Adaptation}: Fine-tune deployed models using small amounts of user-specific data (transfer learning)
    \item \textbf{Diverse Training Data}: Include multi-subject, multi-demographic datasets covering typical variation ranges
\end{enumerate}

\subsubsection{Implications for Real-World Deployment}
These failure modes have practical consequences:
\begin{itemize}
    \item \textbf{Healthcare Monitoring}: Transition confusion could trigger false alerts in fall detection systems or undercount actual activity changes
    \item \textbf{Fitness Tracking}: Walking/Running ambiguity affects calorie estimation and workout classification
    \item \textbf{Rehabilitation}: Gait asymmetry detection requires identifying abnormal patterns, but model trained on healthy subjects may misclassify pathological gaits as normal activities
\end{itemize}

Future framework versions should include:
\begin{enumerate}
    \item \textbf{Confidence Thresholding}: Report prediction probabilities, flagging low-confidence classifications for manual review
    \item \textbf{Uncertainty Quantification}: Bayesian approaches or ensemble disagreement to estimate prediction reliability
    \item \textbf{Online Learning}: Deployed models adapt to individual user patterns, reducing inter-subject variability effects
\end{enumerate}

\subsection{Comparison with Related Work}

\subsubsection{Commercial Platforms}
Our comparison with Edge Impulse and SensiML reveals trade-offs between cost, control, and convenience:

\textbf{Edge Impulse} prioritizes ease-of-use with automated feature engineering and deployment pipelines, but at \$20-100/month subscription costs and limited preprocessing transparency. Our framework matches its accessibility while providing full control and eliminating recurring costs—critical for academic research and resource-limited developing-world contexts.

\textbf{SensiML} targets commercial IoT deployments with sophisticated AutoML capabilities but costs \$99-500/month and supports fewer platforms. Its accuracy on our HAR test (93.1\%) lagged both our framework (96.2\%) and Edge Impulse (94.8\%), possibly due to different algorithm choices or feature sets not optimized for IMU data.

\textbf{Manual TensorFlow Lite} deployment achieved 96.5\% accuracy (marginally better than our 96.2\%), but required expert-level knowledge of model quantization, TFLM runtime integration, and platform-specific optimizations—skills beyond most researchers and hobbyists. Our framework narrows this gap to 0.3\% while requiring no embedded systems expertise.

The conclusion: our framework occupies a unique niche combining professional-grade performance, full transparency, and zero cost—a combination unavailable in existing solutions.

\subsubsection{Academic HAR Systems}
Compared to academic HAR research using similar IMU-based approaches \cite{anguita2013har}, our accuracy (96.2\%) is competitive with published results on benchmark datasets (92-97\% typical). However, most academic systems report only Python prototypes or offline analysis; few provide deployment-ready embedded code. Our contribution lies in bridging this "deployment gap" by automating the translation from trained scikit-learn models to production microcontroller firmware.

Additionally, our systematic treatment of class imbalance addresses a weakness in much HAR literature, where balanced test sets or majority-class-dominated metrics obscure real-world performance issues \cite{he2009learning}.

\subsection{Limitations and Threats to Validity}

\subsubsection{Dataset Limitations}
Our validation uses data from a single subject performing six activities. While sufficient for demonstrating framework capabilities, generalization to diverse populations (different ages, body types, mobility levels) requires multi-subject studies. Activity classes are also limited—real-world applications may need detection of:
\begin{itemize}
    \item Fall events (critical for elderly care)
    \item Complex activities (cooking, cleaning, driving)
    \item Transitions between activities
    \item Abnormal gait patterns (injury, neurological conditions)
\end{itemize}

Future work should validate the framework using publicly available HAR datasets (e.g., UCI HAR, WISDM, PAMAP2) encompassing diverse subjects and extended activity taxonomies.

\subsubsection{Window Size Selection}
The 0.75-second window (75 samples @ 100Hz) is a heuristic choice balancing temporal context and responsiveness. Prior research shows window size significantly impacts accuracy \cite{banos2014wisdm}: shorter windows (<0.5s) may miss full motion cycles, while longer windows (>2s) reduce temporal resolution. Our framework currently uses a fixed window size; adaptive windowing strategies (varying size based on detected activity intensity) could improve performance but add complexity.

\subsubsection{Sensor Placement and Orientation}
Experiments assume the IMU is worn on the wrist in a consistent orientation. Real-world deployment faces variability:
\begin{itemize}
    \item Different body locations (wrist, hip, ankle) produce distinct signal characteristics
    \item Device rotation relative to body axes introduces orientation ambiguity
    \item Multiple simultaneous sensors (e.g., both wrists + hip) require sensor fusion
\end{itemize}

Robust systems require orientation-invariant feature extraction (e.g., computing acceleration magnitude rather than individual axes) or orientation estimation algorithms. Our framework's feature set includes some magnitude-based features, but systematic evaluation of orientation robustness is needed.

\subsubsection{Computational Platform Diversity}
While we demonstrate deployment on ARM Cortex-M4 (Seeed XIAO), the framework claims platform-agnostic code generation. Validation on other architectures—AVR (Arduino Uno), RISC-V (ESP32-C3), ARM Cortex-M0+ (Raspberry Pi Pico)—is needed to confirm portability. Memory and speed characteristics will differ, particularly on 8-bit AVR processors lacking hardware floating-point support.

\subsubsection{External Validity}
Results are specific to human activity recognition using wrist-worn IMU sensors. Applicability to other motion tracking domains (gesture recognition, animal behavior tracking, robotic motion profiling) remains to be validated. The framework's architecture is designed to be domain-agnostic, but preprocessing strategies and feature sets may require adaptation for non-HAR applications.

\subsection{Ethical Considerations}

\subsubsection{Privacy and Data Protection}
Motion tracking systems, even when limited to IMU sensor data, raise important privacy concerns. While accelerometer and gyroscope readings do not capture images or audio, research demonstrates that motion patterns can reveal sensitive information including:
\begin{itemize}
    \item \textbf{Health Status}: Gait abnormalities may indicate neurological conditions, injuries, or degenerative diseases
    \item \textbf{Behavioral Patterns}: Activity logs reveal daily routines, sleep schedules, and lifestyle choices
    \item \textbf{Identity}: Motion signatures can be used for biometric identification, potentially enabling tracking without consent
\end{itemize}

Our framework's edge-first architecture provides inherent privacy advantages: data processing occurs locally on the device, with no requirement for cloud transmission. However, deployment scenarios must consider:
\begin{itemize}
    \item \textbf{Data Storage}: Even local storage of raw sensor data creates risks if devices are lost or stolen
    \item \textbf{Result Transmission}: If activity classifications are transmitted (e.g., via Bluetooth to smartphones), encryption and access controls are essential
    \item \textbf{De-identification}: Research datasets should remove temporal metadata and aggregate statistics to prevent re-identification
\end{itemize}

\subsubsection{Informed Consent}
Deployment of motion tracking systems for research or commercial purposes requires transparent informed consent protocols. Users must understand:
\begin{enumerate}
    \item What data is being collected (raw sensor values vs. derived features vs. activity labels)
    \item How long data is retained and where it is stored
    \item Who has access to the data and for what purposes
    \item Rights to withdraw consent and request data deletion
    \item Potential risks including privacy breaches and unintended uses
\end{enumerate}

The single-subject validation presented in this thesis was conducted under informal self-experimentation protocols appropriate for proof-of-concept research. Future multi-subject studies must obtain institutional review board (IRB) approval and implement formal consent procedures compliant with regulations such as GDPR (Europe) or HIPAA (US healthcare contexts).

\subsubsection{Responsible Use and Potential Misuse}
While the framework was designed for beneficial applications (healthcare monitoring, fitness tracking, accessibility technologies), the same technology could enable harmful uses:
\begin{itemize}
    \item \textbf{Surveillance}: Covert monitoring of employees, students, or vulnerable populations without consent
    \item \textbf{Discrimination}: Using activity patterns to make employment, insurance, or educational decisions
    \item \textbf{Manipulation}: Targeting marketing or interventions based on inferred behavioral patterns
\end{itemize}

As framework developers, we advocate for responsible use guidelines including:
\begin{enumerate}
    \item Mandatory consent mechanisms integrated into deployment templates
    \item Clear documentation discouraging surveillance applications
    \item Open-source licensing (MIT) that preserves user freedoms while disclaiming liability for misuse
    \item Community governance to establish ethical best practices
\end{enumerate}

\subsubsection{Accessibility and Digital Equity}
By eliminating commercial licensing costs (\$20-500/month for competing platforms) and reducing technical barriers, this framework promotes digital equity in edge AI research. However, accessibility challenges remain:
\begin{itemize}
    \item Hardware costs (\$10-30 for development boards) still exclude lowest-resource contexts
    \item Internet access required for initial framework installation and documentation
    \item Technical literacy barriers despite reduced programming requirements
\end{itemize}

Future work should explore partnerships with educational institutions and maker spaces to provide subsidized hardware access and training programs in developing regions.

\subsubsection{Environmental Considerations}
Edge AI systems offer environmental benefits compared to cloud-based approaches by reducing data transmission energy costs and enabling localized processing. Our measured power consumption (30mW average, 92 hours battery life) supports sustainable deployment. However, full lifecycle assessment must consider manufacturing impacts, e-waste from sensor devices, and energy costs of model training infrastructure.

\subsection{Design Trade-offs and Alternatives}

\subsubsection{Classical ML vs Deep Learning}
We deliberately chose classical machine learning (Random Forest, SVM, Neural Networks) over modern deep learning (CNN, LSTM, Transformer) due to:
\begin{itemize}
    \item \textbf{Deployment Constraints}: Classical models compile to small C++ code (95-182KB), while deep networks require specialized inference runtimes (TensorFlow Lite Micro) adding 100-300KB overhead
    \item \textbf{Training Efficiency}: Classical models train in seconds to minutes on CPU; deep networks require GPUs and hours of training
    \item \textbf{Interpretability}: Random Forest feature importance and SVM decision boundaries are auditable; deep networks are opaque
\end{itemize}

However, deep learning excels on large datasets (>100K samples) and raw signal input (skipping manual feature engineering). Future framework versions should integrate TensorFlow Lite support for users with computational resources and large annotated datasets.

\subsubsection{Handcrafted vs Learned Features}
Our 138-feature set uses handcrafted time/frequency domain statistics. This requires domain knowledge (knowing which features discriminate activities) but ensures interpretability and compact representation. End-to-end deep learning could learn features automatically, but at the cost of larger models and reduced transparency.

An intermediate approach—using convolutional autoencoders for learned feature extraction followed by classical ML for classification—could combine benefits. This remains future work.

\subsubsection{Optimization Modes}
The four optimization modes (accuracy, speed, power, balanced) provide flexibility, but require users to understand trade-offs. Alternative designs could:
\begin{itemize}
    \item \textbf{Automated Profiling}: Framework measures device capabilities and selects optimal mode
    \item \textbf{Multi-Objective Optimization}: Pareto frontier exploration to find best accuracy-speed-power compromises
    \item \textbf{Adaptive Runtime}: Device dynamically adjusts inference complexity based on battery level or detected motion intensity
\end{itemize}

Implementing these sophisticated strategies would improve usability but increase framework complexity.

\subsection{Implications for Practice}

\subsubsection{Research and Education}
The framework democratizes edge AI research by eliminating cost barriers (vs \$20-500/month commercial platforms) and technical barriers (vs manual TensorFlow Lite deployment). This enables:
\begin{itemize}
    \item Undergraduate/graduate ML courses to include hands-on edge deployment labs
    \item Researchers in developing countries to conduct HAR studies without institutional budgets for commercial licenses
    \item Rapid prototyping of novel activity recognition algorithms without low-level embedded programming
\end{itemize}

\subsubsection{Healthcare Applications}
Validated models could support wearable health monitoring:
\begin{itemize}
    \item \textbf{Rehabilitation}: Track recovery progress by monitoring activity levels and gait quality
    \item \textbf{Elderly Care}: Detect falls, sedentary behavior, and activity decline
    \item \textbf{Clinical Trials}: Objective activity measurement replacing self-reported logs
\end{itemize}

However, medical deployment requires regulatory compliance (FDA clearance, CE marking) and extensive multi-site validation studies—beyond the scope of this research framework.

\subsubsection{Sports and Fitness}
Consumer wearables (smartwatches, fitness bands) could integrate framework-generated models for:
\begin{itemize}
    \item Automatic workout detection (running, cycling, weightlifting)
    \item Form analysis (detecting improper squat technique, running gait asymmetry)
    \item Personalized coaching (real-time feedback on activity intensity)
\end{itemize}

\subsection{Future Research Directions}

\subsubsection{Short-Term Extensions}
Immediate improvements within the current architecture:
\begin{enumerate}
    \item \textbf{Additional Algorithms}: Integrate XGBoost, LightGBM, k-NN
    \item \textbf{Feature Selection}: Implement automated feature importance analysis to reduce dimensionality
    \item \textbf{Model Compression}: Pruning, quantization-aware training to further reduce memory
    \item \textbf{More Platforms}: ESP32, STM32, nRF5340 deployment support
    \item \textbf{Real-Time Visualization}: Live activity stream from deployed device to framework GUI
\end{enumerate}

\subsubsection{Long-Term Vision}
Transformative enhancements requiring architectural changes:
\begin{enumerate}
    \item \textbf{Deep Learning Integration}: TFLite Micro generation for\\CNN/LSTM models
    \item \textbf{Federated Learning}: Enable collaborative model training across multiple devices without sharing raw data
    \item \textbf{Explainable AI}: Visualize which sensor patterns triggered specific activity predictions
    \item \textbf{Multi-Modal Fusion}: Combine IMU with audio, pressure, or camera data
    \item \textbf{Online Learning}: Deployed models adapt to user-specific patterns over time
    \item \textbf{Cloud Integration}: Optional backend for data backup, model versioning, community model sharing
\end{enumerate}

\subsection{Summary}
This discussion contextualizes our experimental findings within existing literature, acknowledges limitations, and charts future directions. The key takeaway: the proposed framework successfully democratizes edge AI development for motion tracking by combining professional-grade performance (96.2\% accuracy, 12ms inference) with accessibility (zero cost, low learning curve) previously considered mutually exclusive. While limitations exist—single-subject validation, fixed window size, platform diversity—the modular architecture provides a foundation for ongoing enhancement by the research community.